% ---
% RESUMOS
% ---

% resumo em português
\setlength{\absparsep}{18pt} % ajusta o espaçamento dos parágrafos do resumo
\begin{resumo}

Este trabalho apresenta a concepção, implementação e avaliação de um protótipo de sistema \textit{web} voltado ao apoio da preparação do Barema de Atividades Complementares do curso de Ciência da Computação da Universidade Estadual de Santa Cruz (UESC). O problema abordado decorre da execução manual desse processo, que exigia do discente a organização de certificados, o cálculo de cargas horárias, a aplicação de limites regulamentares e a indexação dos comprovantes em um documento único, etapas suscetíveis a erros e retrabalho. A solução proposta centraliza o preenchimento dos dados do estudante, a anexação dos certificados, a pré-validação das informações e a geração de um PDF consolidado, mantendo a análise final sob responsabilidade do colegiado. O protótipo foi desenvolvido com arquitetura em camadas, utilizando Flask, regras de negócio estruturadas, extração textual de documentos, mecanismos complementares de OCR e serviços de geração de PDF. A avaliação foi realizada por meio de cenários funcionais em ambiente local, verificando a aplicação de limites de carga horária, a identificação de inconsistências, a detecção de possíveis duplicidades e a geração do documento final. Os resultados indicam que o sistema contribui para transformar um fluxo manual e fragmentado em um processo assistido, reduzindo falhas formais antes da submissão oficial e oferecendo apoio tanto ao discente quanto à coordenação. Conclui-se que a aplicação construída é adequada como instrumento de organização, pré-validação e padronização do Barema, sem substituir a análise institucional das Atividades Complementares.

 \textbf{Palavras-chave}: Atividades Complementares; Barema; Sistema \textit{web}; Pré-validação.
\end{resumo}


